\section*{Chapitre \ref{ch:b2:field-lab-methods-2} --- Méthodes de terrain et de laboratoire II}
% numbers from the Book 4 second-writer scratch calculator calc/ch31.py

\begin{solution}{exo:b2:field-lab-methods-2:1}
(a) (i) \omterm{def:b2:field-lab-methods-2:textures}{Texture ophitique}~; (ii) \omterm{def:b2:field-lab-methods-2:textures}{texture granoblastique}~; (iii) un \omterm{def:b2:field-lab-methods-2:porphyroblast}{pœciloblaste}, un \omterm{def:b2:field-lab-methods-2:porphyroblast}{porphyroblaste} plein
d'inclusions~; (iv) une \omterm{def:b2:field-lab-methods-2:porphyroblast}{couronne}, une bordure réactionnelle entourant un grenat en cours de décomposition. (b) Par
\eqref{eq:b2:field-lab-methods-2:lami}, $\gamma_1/\gamma_2 = \sin\ang{104}/\sin\ang{128} =
0.9703/0.7880 = 1.23$~: ce joint tire \qty{23}{\percent} plus fort que les deux autres.
\end{solution}

\begin{solution}{exo:b2:field-lab-methods-2:2}
(a) $Q + A + P = 80$~: Q $= 100 \times 25/80 = 31.3$, A $= 18.8$, P $= 50.0$. (b) P/(A+P) $=
40/55 = \qty{72.7}{\percent}$, entre 65 et 90 avec Q compris entre 20 et 60~: une granodiorite.
(c) $M' = 12 + 8 = \qty{20}{\percent}$, en dessous de 35~: leucocrate.
\end{solution}

\begin{solution}{exo:b2:field-lab-methods-2:3}
(a) $\sigma = \sqrt{0.12 \times 0.88/800} = 0.0115$, soit $12.0 \pm 1.1$ points de pourcentage, une
erreur relative de $0.0115/0.12 = \qty{9.6}{\percent}$. (b) Quatre fois plus de points la divisent par deux~: $0.0057$.
\end{solution}

\begin{solution}{exo:b2:field-lab-methods-2:4}
(a) La perte est de \qty{0.0288}{g}, et $0.0288/0.8000 = \qty{3.60}{\percent}$.
(b) $0.7712/8.7712 = 0.0879$.
(c) Poudre calcinée~: $5.20/0.08792 = \qty{59.1}{\percent}$~; poudre d'origine~:
$59.14 \times (1 - 0.0360) = \qty{57.0}{\percent}$.
\end{solution}

\begin{solution}{exo:b2:field-lab-methods-2:5}
(a) $\sigma = 0.10 \times 0.02 = 0.002$ et $n = 0.02 \times 0.98/0.002^2 = 4900$ points.
(b) $\sigma = 0.03$ et $n = 0.30 \times 0.70/0.03^2 = 233.3$, soit 234 points. (c) À erreur relative
fixée, $n = (1-p)/(p\,\varepsilon^2)$ croît comme $1/p$~: le minéral accessoire exige 21 fois plus de
points, davantage qu'une coupe n'en offre en général.
\end{solution}

\begin{solution}{exo:b2:field-lab-methods-2:6}
(a) RGM-1~: $72.89/73.45 = \qty{99.2}{\percent}$~; BCR-2~: $54.17/54.1 = \qty{100.1}{\percent}$.
(b) RGM-1~: $0.19/\sqrt{20} = 0.042$, et $-0.56/0.042 = -13.2$~; BCR-2~: $0.13/\sqrt{20} = 0.029$, et
$0.07/0.029 = 2.4$. (c) RGM-1 présente un biais~: sa moyenne s'écarte de la valeur recommandée de 13 fois son propre
écart type, bien plus que ne le permet le hasard (si la valeur recommandée est juste). BCR-2 est à la
limite de ce que donne le hasard. RGM-1 n'est pas moins précis~: son écart type relatif,
$0.19/72.89 = \qty{0.26}{\percent}$, est proche de celui de BCR-2, \qty{0.24}{\percent}. Le problème est
la justesse, non la précision.
\end{solution}

\begin{solution}{exo:b2:field-lab-methods-2:7}
(a) $M = 39.098 + 3 \times 24.305 + 26.982 + 3 \times 28.085 + 12 \times 15.999 + 2 \times 1.008 =
\qty{417.25}{g.mol^{-1}}$~; les deux OH font une \ce{H2O} ($18.015$), soit \qty{4.32}{\percent}.
(b) \ce{K2O} $47.10/417.25 = \qty{11.3}{\percent}$~; \ce{MgO} $3 \times 40.304/417.25 =
\qty{29.0}{\percent}$~; \ce{Al2O3} $50.98/417.25 = \qty{12.2}{\percent}$~; \ce{SiO2} $3 \times
60.083/417.25 = \qty{43.2}{\percent}$~; total \qty{95.7}{\percent}. (c) La microsonde ne détecte pas l'hydrogène~: une analyse de mica
manque son eau, environ \qty{4}{\percent}, et son total doit être voisin de
96.
\end{solution}

\begin{solution}{exo:b2:field-lab-methods-2:8}
(a) $\beta = \ln(0.1194/0.11905)/\ln(85.909261/87.905612) = 0.002936/(-0.02297) = -0.1278$.
(b) $0.71062 \times (86.908878/85.909261)^{-0.1278} = 0.70957$. (c) Le 86/88 mesuré est inférieur au
rapport vrai~: l'isotope lourd a été favorisé, et le 87/86 mesuré est trop élevé.
\end{solution}

\begin{solution}{exo:b2:field-lab-methods-2:9}
(a) $0.1000 \times 15.271 \times 10^{-6}/83.9242 \times 1.3599 \times 10^{-3} = \qty{2.474e-11}{mol}$.
(b) Par \eqref{eq:b2:field-lab-methods-2:id}, $N_x = 2.474 \times 10^{-11} \times (732.8 -
0.874)/(0.874 - 0.0568) = \qty{2.216e-8}{mol}$ de $^{86}\mathrm{Sr}$~; strontium~: $2.216 \times
10^{-8}/0.0986 \times 87.62 = \qty{1.969e-5}{g}$ dans \qty{0.0500}{g}, soit \qty{394}{ppm} (393.9).
(c) Un recouvrement de \qty{99.95}{\percent}~: la méthode est exacte à l'arrondi du rapport près.
\end{solution}

\begin{solution}{exo:b2:field-lab-methods-2:10}
(a) $R_m = \sqrt{732.8 \times 0.05680} = 6.45$, où $\mathcal M = 1.018$. (b) $\mathcal M = 2.31$
à 0.1, 1.13 à 0.5, 1.038 à 1.664 et 1.16 à 100. (c) Près de $R_x$ le mélange ne diffère guère de
l'échantillon, près de $R_s$ du spike, et la différence qui porte l'information se perd dans
l'erreur de $R_m$~: $\mathcal M$ diverge aux deux extrémités. Entre elles, le minimum est très plat
($\mathcal M < 1.2$ de 0.4 environ à 100), si bien que l'analyste n'a besoin que du bon ordre de grandeur
du spike, qu'une concentration grossière, par exemple par XRF, fournit.
\end{solution}

\begin{solution}{exo:b2:field-lab-methods-2:11}
(a) $\sigma = \sqrt{0.62 \times 0.38/400} = 0.0243$~: $62 \pm 2.4$. (b) $(65 - 62)/2.43 = 1.2$
écarts types~: la roche peut bien être une granodiorite. (c) $\sigma = 0.015$ exige $0.62 \times
0.38/0.015^2 = 1047.1$, soit 1048 points de feldspath.
\end{solution}

\begin{solution}{exo:b2:field-lab-methods-2:12}
(a) Chaque point manque avec la probabilité $1 - p$, indépendamment~: $(1-p)^n$. (b) $0.99^n \le 0.05$
donne $n \ge \ln 0.05/\ln 0.99 = 298.1$~: 299 points. (c) $n \ge \ln 0.05/\ln 0.9995 = 5990$~: un comptage
sur grille est le mauvais outil~; un minéral accessoire rare se recherche en balayant toute la section.
\end{solution}

\begin{solution}{pb:b2:field-lab-methods-2:1}
\noindent\textbf{Partie I.}
\textbf{1.} Quartz $372/1200 = \qty{31.0}{\percent}$, feldspath alcalin 29.0, plagioclase 27.0, biotite
9.0, muscovite 3.0, oxydes opaques 1.0.
\textbf{2.} $\sqrt{0.31 \times 0.69/1200} = 0.0134$~; muscovite~: $\sqrt{0.97/(0.03 \times 1200)} =
\qty{16.4}{\percent}$.
\textbf{3.} $M = 9 + 3 + 1 = \qty{13.0}{\percent}$~; $M' = 13 - 3 = \qty{10.0}{\percent}$ sans la
muscovite~: leucocrate.
\textbf{4.} $Q + A + P = 1044$ points~: Q $= 35.6$, A $= 33.3$, P $= 31.0$.
\textbf{5.} $324/672 = \qty{48.2}{\percent}$, avec Q compris entre 20 et 60~: un granite, et entre 35 et
65~: un monzogranite.
\textbf{6.} $\sqrt{0.482 \times 0.518/672} = 0.0193$~: les seuils à 35 et 65 sont à 6.9 et 8.7
écarts types. Le nom est sûr.
\textbf{7.} $0.31 \times 0.69/0.01^2 = 2139$ points.

\noindent\textbf{Partie II.}
\textbf{8.} $0.0051/0.8000 = \qty{0.64}{\percent}$.
\textbf{9.} $0.7949/8.7949 = 0.0904$, et $72 \times 0.0904 = \qty{6.51}{\percent}$ \ce{SiO2} dans le
verre~: la fusion dilue l'échantillon environ onze fois.
\textbf{10.} $\bar I = 24.91$, $\bar c = 62.283$, $\sum(I - \bar I)^2 = 31.918$, $\sum(I - \bar I)(c -
\bar c) = 80.009$~: $a = 2.5067$ par kcps, $b = 62.283 - 2.5067 \times 24.91 = -0.16$.
\textbf{11.} $2.5067 \times 28.63 - 0.16 = \qty{71.61}{\percent}$.
\textbf{12.} BCR-2 $-0.04$, AGV-2 $+0.05$, RGM-1 $-0.01$~: toutes inférieures aux écarts types
de la méthode (0.1 à 0.2), donc la droite est adéquate sur cet intervalle.
\textbf{13.} $71.61 + 27.15 + 0.64 = \qty{99.40}{\percent}$, entre 97 et 101~: non reprise.
\textbf{14.} $71.61 \times 100/(71.61 + 27.15) = \qty{72.51}{\percent}$.

\noindent\textbf{Partie III.}
\textbf{15.} $0.0986/0.8258 = 0.1194$~: la valeur de normalisation est le rapport naturel.
\textbf{16.} $\beta = 0.007313/(-0.022972) = -0.3183$.
\textbf{17.} $0.71655 \times (86.908878/85.909261)^{-0.3183} = 0.71392$.
\textbf{18.} Le 86/88 mesuré vaut $0.11853/0.1194 - 1 = \qty{-0.73}{\percent}$ par rapport à la valeur vraie
sur deux unités de masse~: environ \qty{0.36}{\percent} par unité de masse, en faveur des isotopes lourds.
\textbf{19.} Le $^{86}\mathrm{Sr}/^{88}\mathrm{Sr}$ du spike est 0.7661~: le 86/88 du mélange n'est plus le
naturel 0.1194, et le fixer à 0.1194 traiterait le spike comme un biais instrumental.

\noindent\textbf{Partie IV.}
\textbf{20.} Somme des rapports à 88~: $561.4 + 0.7661 + 0.1819 + 1 = 563.348$~; $561.4/563.348 =
\qty{99.65}{\percent}$ et $0.7661/563.348 = \qty{0.13599}{\percent}$.
\textbf{21.} $0.2500 \times 15.271 \times 10^{-6}/83.9242 = \qty{4.549e-8}{mol}$ de strontium, et
$\times 1.3599 \times 10^{-3} = \qty{6.186e-11}{mol}$ de $^{86}\mathrm{Sr}$.
\textbf{22.} $561.4/0.7661 = 732.8$ et $0.0056/0.0986 = 0.05680$.
\textbf{23.} $N_x = 6.186 \times 10^{-11} \times (732.8 - 1.6812)/(1.6812 - 0.0568) =
\qty{2.784e-8}{mol}$.
\textbf{24.} $\mathcal M = 1.6812 \times 732.74/(731.12 \times 1.6244) = 1.04$~; le rapport ajoute
$1.04 \times \qty{0.01}{\percent}$, le spike $0.031/15.271 = \qty{0.20}{\percent}$~; en quadrature
\qty{0.20}{\percent}, dominée par l'étalonnage du spike.
\textbf{25.} Strontium~: $2.784 \times 10^{-8}/0.0986 = \qty{2.824e-7}{mol}$, soit $2.824 \times
10^{-7} \times 87.62 = \qty{2.474e-5}{g}$ dans \qty{0.1000}{g}. Le granite de Varden contient
\textbf{\qty[text-series-to-math=true]{247}{ppm}} de strontium, $\pm 0.5$.
\end{solution}
